\section{Introduction}

\subsection{Application setup}

This work investigates the numerical properties of sparse linear
systems arising from the semi-implicit discretization of the
two-dimensional linearized shallow water equations. The shallow
water model is a depth-averaged approximation of fluid dynamics
that is applicable when the horizontal length scales are much
larger than the vertical depth. It describes the evolution of
the free-surface elevation and the depth-averaged horizontal
velocity, making it a useful model for studying wave propagation
in coastal and oceanographic settings
\cite{casulli1990,chopra2025}.

Let $\Omega\subset\mathbb{R}^2$ be a rectangular computational
domain, and let $t\geq 0$ denote time. The variables
$\eta(x,y,t)$ and $\mathbf{u}(x,y,t)=(u,v)^\top$ represent,
respectively, the free-surface displacement relative to a resting
state and the depth-averaged horizontal velocity. The function
$H(x,y)>0$ denotes the undisturbed water depth, $g$ is the
gravitational acceleration, and $f$ is a constant Coriolis
parameter.

Under a linearization about a resting state, neglecting nonlinear
advection and friction, the governing equations are
\begin{align}
    \frac{\partial \eta}{\partial t}
    + \nabla\cdot(H\mathbf{u})    & = 0,
    \label{eq:swe_continuity}
    \\
    \frac{\partial \mathbf{u}}{\partial t}
    + g\nabla\eta
    + f\mathbf{k}\times\mathbf{u} & = 0,
    \label{eq:swe_momentum}
\end{align}
where $\mathbf{k}$ denotes the vertical unit vector.

The first equation expresses conservation of mass, while the
second describes the linearized horizontal momentum balance.
The simplified model isolates the gravity-wave dynamics and
provides a controlled setting in which the algebraic properties
of the numerical discretization can be studied independently
of a full nonlinear oceanographic simulation.

The initial experiments will use a periodic rectangular domain.
A constant reference depth, $H(x,y)=H_0$, will provide the
baseline configuration. Subsequently, spatially heterogeneous
depth profiles will be introduced while maintaining
$H(x,y)>0$. Recent studies on shallow water models and
structure-preserving numerical formulations provide additional
context for these modeling and discretization choices
\cite{cardosoribeiro2024,chopra2025}.

The spatial discretization will be based on finite differences
in conservative flux form. For the temporal discretization, the
gravity-wave terms will be treated implicitly, while the Coriolis
term will initially be treated explicitly. This semi-implicit
strategy couples the unknown free-surface elevation at the new
time level to the unknown velocity at that same level.

Writing $\Delta t$ for the time step and using superscripts
$n$ and $n+1$ for consecutive time levels, the discrete equations
are
\begin{align}
    \frac{\eta^{n+1}-\eta^n}{\Delta t}
    + \nabla_h\cdot(H\mathbf{u}^{n+1}) & = 0,
    \label{eq:discrete_continuity}
    \\
    \frac{\mathbf{u}^{n+1}-\mathbf{u}^n}{\Delta t}
    + g\nabla_h\eta^{n+1}
    + f\mathbf{k}\times\mathbf{u}^{n}  & = 0,
    \label{eq:discrete_momentum}
\end{align}
where $\nabla_h$ and $\nabla_h\cdot$ denote the discrete gradient
and divergence operators.

Eliminating the new velocity from these equations leads to a
linear system for the free-surface elevation:
\begin{equation}
    A_h\eta^{n+1}=b_h^n,
    \label{eq:linear_system}
\end{equation}
where
\begin{equation}
    A_h =
    I-g\Delta t^2
    \nabla_h\cdot\bigl(H\nabla_h\bigr).
    \label{eq:discrete_operator}
\end{equation}
Here, $I$ is the identity operator and $b_h^n$ is determined by
the solution at the previous time level.

The matrix representation of $A_h$ is sparse because the
finite-difference operators couple each grid point only to a
small number of neighboring points. This sparsity makes the
problem suitable for iterative linear solvers.

A compatible flux-form discretization will be selected so that
the resulting matrix is symmetric positive definite under the
adopted boundary conditions. In particular, if the discrete
divergence and gradient operators are negative adjoints and
$H$ is represented by a positive diagonal matrix $H_h$, then
\begin{equation}
    A_h =
    I+g\Delta t^2 G_h^\top H_hG_h,
    \label{eq:spd_operator}
\end{equation}
where $G_h$ is the discrete gradient matrix. This representation
provides a mathematical basis for investigating the Conjugate
Gradient method. Symmetry and positive definiteness will be
verified for the implemented system before this method is used.

The classical semi-implicit shallow water formulation of
Casulli \cite{casulli1990} and the conjugate-gradient approach
studied by Bonaventura and Rosatti \cite{bonaventura2002}
provide the principal methodological foundations for this
part of the work.

The experiments will investigate two main factors: spatial
resolution and coefficient heterogeneity. First, the grid will be
progressively refined while the time step and the remaining
parameters are held fixed. This experiment will examine how
the dimension and spectral properties of $A_h$
change with spatial resolution.

Second, the depth coefficient will be varied while the grid and
time step remain fixed. The heterogeneity will be quantified by
the contrast ratio
\begin{equation}
    \chi=\frac{H_{\max}}{H_{\min}},
\end{equation}
where $H_{\min}>0$. This experiment will investigate whether
changes in coefficient contrast are associated with changes
in the conditioning of the matrix and the convergence behavior
of the iterative solvers.

The principal methods to be compared are the unpreconditioned
Conjugate Gradient method, Jacobi-preconditioned Conjugate
Gradient, and incomplete-Cholesky-preconditioned Conjugate
Gradient. The choice of these methods is motivated by the
literature on iterative methods for sparse linear systems and
preconditioning techniques \cite{saad2003,benzi2002}. More
recent research also investigates advanced preconditioning
strategies for hyperbolic PDE systems, including numerical
experiments involving shallow water equations
\cite{krzysik2025}.

The evaluation will record the iteration count, relative residual
norm, estimated condition number, and computational time.
Preconditioner construction time and solution time will be
reported separately. Since the coefficient matrix remains
unchanged across time steps when the grid, depth profile, and
time step are fixed, the study will also examine the effect of
reusing a preconditioner across repeated linear solves.

This cost-aware comparison constitutes the principal
methodological extension of the work: rather than evaluating
preconditioners only by their iteration counts, it will examine
when the reduction in solution cost compensates for the
additional setup cost. The implementation will use consistent
stopping tolerances and initial guesses to ensure meaningful
comparisons.

The main research question is therefore:

\begin{quote}
    How do spatial refinement and depth-coefficient heterogeneity
    affect the conditioning and iterative convergence of sparse
    linear systems arising from the semi-implicit shallow water
    equations, and under what conditions does incomplete Cholesky
    preconditioning provide a computational advantage over
    unpreconditioned and Jacobi-preconditioned Conjugate Gradient?
\end{quote}